\documentclass[sigconf]{acmart}

\settopmatter{printacmref=false}
\renewcommand\footnotetextcopyrightpermission[1]{}
\setcopyright{none}

\acmConference[SBES 2026]{40th Brazilian Symposium on Software Engineering}{September 8--12, 2026}{São Paulo, SP, Brazil}

\renewcommand\abstractname{Resumo}
\renewcommand\keywordsname{Palavras-chave}
\renewcommand\refname{Referências}

\AtBeginDocument{
  \providecommand\BibTeX{{Bib\TeX}}
}

\usepackage{booktabs}
\usepackage{listings}
\usepackage{xcolor}
\usepackage{url}
\usepackage{mdframed}
\usepackage{tikz}
\usepackage{iftex}
\iftutex
  \usepackage{fontspec}
\else
  \usepackage[utf8]{inputenc}
  \usepackage[T1]{fontenc}
\fi
\usepackage[brazil]{babel}
\usepackage{hyperref}
\usetikzlibrary{calc}

\newmdenv[
  leftline=true,
  rightline=false,
  topline=false,
  bottomline=false,
  linewidth=3pt,
  linecolor=blue!50!black,
  backgroundcolor=blue!3!white,
  innerleftmargin=8pt,
  innerrightmargin=8pt,
  innertopmargin=4pt,
  innerbottommargin=4pt,
  skipabove=6pt,
  skipbelow=6pt
]{promptbox}

\begin{document}

\title{Dívida Técnica Arquitetural em Sistemas Gerados por IA Generativa: um estudo experimental sobre MVPs e escalabilidade}

\author{Anne Esther Lins Figueirôa}
\affiliation{
  \institution{Instituto de Tecnologia e Liderança -- INTELI}
  \city{São Paulo}
  \country{Brazil}
}
\email{anne.figueiroa@inteli.edu.br}

\author{Reginaldo Arakaki}
\affiliation{
  \institution{Instituto de Tecnologia e Liderança -- INTELI}
  \city{São Paulo}
  \country{Brazil}
}
\email{reginaldo.arakaki@inteli.edu.br}

\author{Fabiana Martins de Oliveira}
\affiliation{
  \institution{Universidade de São Paulo -- USP}
  \city{São Paulo}
  \country{Brazil}
}
\email{fabiana.martins@usp.br}

\renewcommand{\shortauthors}{Figueirôa et al.}
\renewcommand{\shorttitle}{Dívida Técnica Arquitetural em Soluções Geradas por IA}

\begin{abstract}
Ferramentas de Inteligência Artificial Generativa (IAG) aceleraram o desenvolvimento
de produtos em startups (empresas de base tecnológica em estágio inicial), permitindo
a construção de Produtos Mínimos Viáveis (MVPs) em ciclos extremamente curtos.
Contudo, pouco se sabe sobre a qualidade arquitetural dos sistemas assim gerados
quando submetidos a condições reais de escala. Este artigo apresenta um estudo
experimental controlado que avalia a dívida técnica arquitetural introduzida por
ferramentas de IAG em comparação com arquiteturas otimizadas via prompts
(instruções em linguagem natural fornecidas a modelos de IA) arquiteturais explícitos.
A partir de uma persona (perfil sintético representativo) baseada em entrevista com
fundadores de uma startup real, implementou-se um backend (camada de servidor
responsável pela lógica de negócio) em dois cenários: Cenário~A, gerado com prompts
de negócio sem especificação arquitetural; e Cenário~B, refatorado com quatro decisões
arquiteturais documentadas. Os cenários foram submetidos a testes de carga com a
ferramentas de teste de carga. Os resultados demonstraram que o Cenário~A colapsou em 2.413 usuários
virtuais simultâneos (24\% da capacidade-alvo de 5.000), com 61,94\% de taxa de
erro e latência superior a 5.000\,ms, evidenciando dívida técnica arquitetural
inadvertida introduzida por prompts sem especificação de engenharia. O Cenário~B,
refatorado com quatro decisões arquiteturais explícitas, completou o protocolo até
10.000 usuários virtuais simultâneos (o dobro da meta), sem nenhum erro e com
latência média de 34\,ms. Como contribuição prática, propõe-se um modelo de prompt arquitetural
de quatro componentes que permite a founders (fundadores sem formação técnica)
especificar restrições de escala de forma que os modelos de IAG consigam traduzir em
padrões de implementação corretos.
\end{abstract}

\renewcommand\abstractname{Abstract}
\begin{abstract} Artificial Intelligence (GenAI) tools have accelerated product development
in startups, enabling the construction of Minimum Viable Products (MVPs) in extremely
short cycles. However, little is known about the architectural quality of systems thus
generated when subjected to real-world scale conditions. This paper presents a
controlled experimental study that evaluates the architectural technical debt introduced
by GenAI tools compared to architectures optimized via explicit architectural prompts.
Based on a persona derived from interviews with founders of a real startup, a backend
was implemented in two scenarios: Scenario A, generated with business prompts without
architectural specification; and Scenario B, refactored with four documented
architectural decisions. Both scenarios were submitted to load testing with dedicated load testing tools.
Results showed that Scenario A collapsed at 2,413 simultaneous virtual users
(24\% of the 5,000-user target), with a 61.94\% error rate and latency exceeding
5,000\,ms, evidencing inadvertent architectural technical debt introduced by prompts
lacking engineering specification. Scenario B, refactored with four explicit
architectural decisions, completed the protocol up to 10,000 simultaneous virtual
users (twice the target), with zero errors and a mean latency of 34\,ms.
As a practical contribution, a four-component architectural prompt model is proposed, enabling
founders without technical background to specify scalability constraints in a way that
GenAI models can translate into correct implementation patterns.
\end{abstract}

\keywords{Dívida Técnica Arquitetural, Inteligência Artificial Generativa, Escalabilidade,
  Prompt Engineering, Startup, Teste de Carga, Decisão Arquitetural, Táticas Arquiteturais}

\frenchspacing
\maketitle
\sloppy

%% ─────────────────────────────────────────────────────────────────────────────
\section{Introdução}
%% ─────────────────────────────────────────────────────────────────────────────

Em 2024, 60\% das startups utilizavam ferramentas \textit{no-code}/\textit{low-code}
e 74\% dos empreendedores incorporavam IA como componente central de seus
produtos~\cite{RudysAI2024,HubSpot2024}. Ferramentas como Lovable, v0 e Replit Agent
permitem construir MVPs em horas via prompts, consolidando o \textit{vibe coding}: desenvolvimento
guiado inteiramente por descrições de negócio, sem especificação de requisitos
técnicos~\cite{HubSpot2024}. Embora eficazes na validação funcional, essas ferramentas
abstraem decisões arquitetônicas críticas --- especialmente de escalabilidade, o
atributo mais negligenciado por IAG~\cite{Nascimento2023,Bass2021}, que se manifesta
em decisões interdependentes e de difícil reversão desde as fases iniciais.

A história confirma os riscos: Twitter colapsou ao escalar seu monólito; Uber migrou
dolorosamente para microsserviços~\cite{Newman2021}. Pesquisas recentes focam correção
sintática de código gerado por IA~\cite{Perry2023,Liu2023,Shypula2023}, mas há
escassez de estudos sobre \textit{integridade arquitetural sistêmica}.

\textbf{Hipótese central:} a qualidade arquitetural do código gerado por IAG depende
diretamente do nível de especificação arquitetural nos prompts --- prompts com
restrições técnicas explícitas produzem sistemas escaláveis; prompts de negócio
sem tal especificação produzem dívida técnica arquitetural inadvertida.

\textbf{Objetivo geral:} avaliar experimentalmente a escalabilidade de arquiteturas
geradas por IAG versus arquiteturas otimizadas via prompts arquiteturais explícitos,
sobre a mesma stack (Node.js/PostgreSQL). Os objetivos específicos são:
(I)~mapear táticas de escalabilidade~\cite{Bass2021} e refatoração~\cite{Pressman2021};
(II)~implementar o \textbf{Cenário~A} via Lovable com prompts de negócio sem especificação
arquitetural; (III)~implementar o \textbf{Cenário~B} com quatro táticas explícitas
(JOIN, pool otimizado, \textit{cache} e fila assíncrona de escritas);
(IV)~submeter ambos a testes de carga~\cite{grafana2024}; e
(V)~comparar ponto de ruptura e eficiência via ATAM~\cite{Kazman2000}.

\section{Fundamentação teórica}
%% ─────────────────────────────────────────────────────────────────────────────

\subsection{Dívida técnica e sua dimensão arquitetural}

O conceito de Dívida Técnica foi cunhado por Cunningham~\cite{Cunningham1992} e
sistematizado por Fowler~\cite{Fowler2018} no ``Quadrante da Dívida Técnica'',
classificando-a em prudente/imprudente e deliberada/inadvertida. Em startups que
adotam IAG predomina a dívida \textit{imprudente-inadvertida}: a equipe não percebe
o passivo por falta de repertório técnico. Fernández-Sánchez et
al.~\cite{FernandezSanchez2020} identificam que a DTA é a de maior impacto no ciclo
de vida dos sistemas --- sistêmica, não localizada --- tendo na ISO/IEC
25010~\cite{ISO25010} seu referencial normativo de mensuração.

\subsection{Escalabilidade como atributo arquitetural}

Bass, Clements e Kazman~\cite{Bass2021} distinguem escalabilidade horizontal
(adição de instâncias) e vertical (aumento de recursos). Sistemas gerados sem
especificação arquitetural tendem a incorporar estado compartilhado em memória e
chamadas síncronas em cascata, que inviabilizam escalonamento horizontal.
Kleppmann~\cite{Kleppmann2017} descreve os tradeoffs centrais (CAP, latência vs.\
throughput) e confirma que a menor complexidade imediata (consultas síncronas, sem
cache, processamento sequencial) corresponde ao maior custo sob escala. As táticas
aplicadas neste estudo~\cite{Bass2021} são: filas assíncronas, \textit{caching}
distribuído e separação de modelos de leitura e escrita.

\subsection{Estado da arte e trabalhos relacionados}

A literatura sobre qualidade de código gerado por IAG concentra-se em dois eixos,
ambos distintos da contribuição central deste trabalho. O primeiro aborda
\textbf{correção funcional}: Perry et al.~\cite{Perry2023} demonstraram que a
qualidade do código depende criticamente do nível de especificação no prompt; Liu et
al.~\cite{Liu2023} identificaram degradação em problemas com múltiplos componentes
interdependentes; e Shypula et al.~\cite{Shypula2023} introduziram o conceito de
``lacuna de desempenho sistêmico'' para descrever a diferença entre código
funcionalmente correto e código eficiente sob operação real. Nenhum desses trabalhos
avalia a integridade arquitetural de sistemas completos. O segundo eixo aborda
\textbf{requisitos não-funcionais}: Nascimento et al.~\cite{Nascimento2023}
identificaram lacunas sistemáticas na geração desses requisitos por modelos
generativos; Vogelsang~\cite{Vogelsang2024} argumenta que requisitos precisam ser
enriquecidos com decisões de design antes de serem usados como prompts,
enriquecimento que denomina trabalho de RE fundamental irredutível à automação. Em
síntese, LLMs\footnote{Modelos de Linguagem de Grande Escala: sistemas de IA
treinados em grandes volumes de texto para gerar linguagem e código.} produzem código
funcionalmente plausível, mas negligenciam sistematicamente atributos de qualidade
não-funcionais. O presente trabalho avança ao avaliar a DTA no nível sistêmico, via testes de carga em condições controladas, e ao propor um modelo operacional de
prompt arquitetural acessível a founders sem formação técnica.

%% ─────────────────────────────────────────────────────────────────────────────
\section{Metodologia}
%% ─────────────────────────────────────────────────────────────────────────────

Esta IC foi desenvolvida individualmente no INTELI (orientador: Prof.\ Arakaki;
coorientadora: Prof.\textsuperscript{a} Fabiana Oliveira/USP). Adota Pesquisa
Experimental Aplicada~\cite{Wohlin2012} em quatro etapas: construção de persona,
implementação dos cenários, testes de carga e análise comparativa. Artefatos em
repositório público\footnote{\url{https://github.com/anneestherlf/welcome-navigator}}~\cite{WelcomeNavigator2025}.

\textbf{Coleta de dados.} Entrevista semiestruturada (30--45\,min) com dois
fundadores de startup pré-seed cujo produto foi construído integralmente com IAG
sem engenheiros~\cite{Ries2012}. Participantes pseudonimizados (TCLE assinado).
Dados analisados pela análise temática de Braun e Clarke~\cite{Braun2006}, gerando
a persona \textit{Onboard}: plataforma SaaS multi-tenant de onboarding corporativo,
fundada sem formação em ES, construída com Lovable e Supabase~\cite{Supabase2024}.
Cinco padrões emergentes: \textbf{P1}~(prompts sem parâmetros arquiteturais),
\textbf{P2}~(adoção acrítica de ferramentas), \textbf{P3}~(aceitação acrítica de
outputs), \textbf{P4}~(arquitetura do protótipo em produção) e \textbf{P5}~(reconhecimento
retrospectivo da lacuna de ES).

\textbf{Cenários experimentais.} O experimento implementa o módulo de trilhas de
aprendizado da Onboard, componente de maior carga computacional, representativo do
padrão Fan-out~\cite{Kleppmann2017} (múltiplos usuários acessando conteúdo
personalizado simultaneamente), sobre a mesma stack Node.js/PostgreSQL.
O \textbf{Cenário~A} é gerado integralmente pelo Lovable a partir de prompts de
negócio sem qualquer especificação arquitetural; as decisões de acesso a dados,
gerenciamento de conexões e estratégia de escritas são delegadas à IAG, espelhando
os padrões P1--P3 da persona.
O \textbf{Cenário~B} reconstrói o mesmo backend via prompts arquiteturais enviados
ao Claude (claude-sonnet-4-5), com quatro instruções explícitas que seguem o modelo
de quatro componentes proposto na Seção~\ref{subsec:modelo_prompt}: (i)~JOIN
eliminando o padrão N+1; (ii)~pool otimizado (\texttt{max:20}, timeout 2\,s,
fail-fast 503); (iii)~\textit{cache} em memória com TTL 60--300\,s; e
(iv)~fila assíncrona de escritas. A stack é idêntica em ambos, de modo que qualquer
diferença observada é atribuível exclusivamente às decisões arquiteturais.

\textbf{Avaliação e protocolo.} Ambos são avaliados pelo
ATAM~\cite{Kazman2000} com SLA definido em 5.000 VUs simultâneos, latência
$<$\,500\,ms e taxa de erro $<$\,1\%. Os testes de carga~\cite{grafana2024}
simulam o padrão Fan-out: cada usuário virtual realiza, em sequência, autenticação,
acesso à trilha e registro de progresso. A carga é injetada em cinco estágios de
60\,s cada: E1 (100 VUs), E2 (500), E3 (1.000), E4 (5.000) e E5 (10.000). Os
cenários rodam em contêineres Docker com 2\,vCPUs e 4\,GB de RAM, sem escalonamento
automático, garantindo que qualquer diferença seja atribuída à arquitetura, não à
infraestrutura. O protocolo é interrompido ao superar 10\% de erros por 15\,s ou
latência $>$\,5.000\,ms. Três métricas estáticas complementam a análise dinâmica:
Complexidade Ciclomática (CC)~\cite{McCabe1976}, Acoplamento entre Objetos (CBO)
e Falta de Coesão em Métodos (LCOM)~\cite{Chidamber1994}.

\section{Cenário A: MVP gerado por IAG}
%% ─────────────────────────────────────────────────────────────────────────────

O Cenário~A foi gerado pelo Lovable via prompt descritivo: ``Crie uma plataforma
de onboarding com banco de dados, login, trilhas de aprendizado, módulos com texto,
barra de progresso e painel do gestor'', sem qualquer parâmetro arquitetural.
Um achado metodologicamente relevante: a resposta ``Admin + Gestor + Colaborador''
às perguntas de clarificação da ferramenta gerou automaticamente RBAC
(\textit{Role-Based Access Control}) completo com RLS (\textit{Row-Level Security})
em 9 tabelas e três \textit{security definer functions}, sem instrução explícita,
evidenciando elicitação implícita de requisitos funcionais pela IAG --- com qualidade
arquitetural variável nos demais viewpoints~\cite{ISO25010}.

\textbf{Pontos de sensibilidade identificados pelo ATAM.} O dashboard do gestor
(\texttt{app.team.tsx}) executa $3+2M$ queries por acesso ($M$ = colaboradores):
com $M=50$, uma única visita dispara 103 queries em loop \texttt{for/await}
síncrono. O total identificado é de 49 chamadas em 5 arquivos (média: 9,8
queries/rota). Pool de conexões com \texttt{max:10} sem timeout satura com
$\approx$10 requisições paralelas. Todas as escritas são síncronas e bloqueantes.
A análise ATAM identificou três tradeoffs: N+1 vs.\ simplicidade de código;
ausência de \textit{cache} vs.\ consistência de dados; e escritas síncronas
vs.\ throughput sob pico.

\section{Cenário B: Arquitetura refatorada}
%% ─────────────────────────────────────────────────────────────────────────────

\subsection{O modelo de prompt arquitetural para founders}
\label{subsec:modelo_prompt}

O Cenário~B foi construído aplicando quatro refatorações ao backend do Cenário~A,
cada uma via prompt explícito ao Claude (claude-sonnet-4-5). Vogelsang~\cite{Vogelsang2024}
argumenta que requisitos precisam ser enriquecidos com decisões de design antes de
serem usados como prompts --- trabalho de RE irredutível à automação. O modelo
proposto operacionaliza essa decomposição em linguagem acessível a founders, sem
formação técnica, como um \textit{architectural constraint prompt pattern}~\cite{White2023}.
Founders identificam funcionalidades lentas pelos sinais: \textbf{S1}~(lentidão
progressiva, indicando N+1), \textbf{S2}~(timeouts sob pico, indicando pool saturado)
e \textbf{S3}~(dashboard lento, indicando ausência de \textit{cache}), e enviam
o prompt com o código existente à IAG. O modelo é projetado exclusivamente para
\textbf{avaliar e corrigir escalabilidade de funcionalidades existentes} --- não
para construir novas. A Tabela~\ref{tab:modelo_prompt} apresenta os quatro
componentes; template completo disponível no repositório~\cite{WelcomeNavigator2025}.

\begin{table}[t]
  \caption{Template de prompt arquitetural: quatro componentes replicáveis.}
  \label{tab:modelo_prompt}
  \small
  \begin{tabular}{p{1.1cm}p{2.6cm}p{2.9cm}}
    \toprule
    \textbf{Comp.} & \textbf{O que escrever} & \textbf{Exemplo} \\
    \midrule
    (1) Função & Descreva a tela/endpoint lento & ``Painel do gestor com lista de colaboradores'' \\
    (2) Restrição & Como o banco deve ser acessado & ``Use JOIN, não loops'' \\
    (3) Tradeoff & Consequência aceita & ``Aceito dados com 30\,s de atraso'' \\
    (4) Docs & Peça explicação do padrão corrigido & ``Comente o que mudaria para 10$\times$ mais usuários'' \\
    \bottomrule
  \end{tabular}
\end{table}

\subsection{As quatro refatorações e seus prompts}

\textbf{R1 — Eliminação do N+1 (S1, S3).} \texttt{GET /app/team} executava um
loop \texttt{for/await} com 2 queries síncronas por colaborador, totalizando 103
queries/req com $M=50$. Restrição central [comp.~2]: \textit{``Use uma única query
SQL com JOINs e GROUP BY --- não execute consultas dentro de loops.''}
Resultado: 103→1 query/req (−99\%), via LEFT JOINs com
\texttt{COUNT(DISTINCT) FILTER (WHERE)}, independente de $M$.

\textbf{R2 — Pool otimizado (S2).} Pool com \texttt{max:10} sem timeout esgotava
em $\approx$10 requisições paralelas. Restrição central [comp.~3]: \textit{``Aceito
que com pool esgotado o sistema retorne 503 imediatamente em vez de requisições
suspensas.''} Resultado: pool 10→20, \textit{fail-fast} com 503 imediato,
eliminando requisições penduradas indefinidamente.

\textbf{R3 — Cache em memória (S3).} Endpoints quase-estáticos (\texttt{GET /app/tracks},
\texttt{GET /app/modules/:id}) iam ao banco a cada requisição. Com 5.000 VUs, cada
segundo gerava milhares de queries idênticas. Restrição central [comp.~3]:
\textit{``Aceito que trilhas editadas demorem 60\,s para aparecer atualizadas.''}
Resultado: $>$98\% de \textit{cache hit} esperado para trilhas e $>$99\% para
módulos a 5.000 VUs.

\textbf{R4 — Fila assíncrona (S2).} \texttt{POST /app/modules/:id/complete}
executava \textit{upsert} síncrono, bloqueando HTTP até a escrita completar.
Restrição central [comp.~3]: \textit{``Aceito perda de progresso em crash do
servidor.''} Resultado: endpoint retorna imediatamente (\texttt{\{queued:true\}});
\textit{worker} processa em 500\,ms; consistência eventual aceitável para onboarding.

As quatro refatorações atuam em camadas complementares sobre o mesmo gargalo:
R1 elimina a causa primária do N+1; R2 garante que o ganho de R1 se traduza em
capacidade real; R3 reduz carga residual nos endpoints de leitura; R4 libera o
pool das escritas síncronas. Dois pontos residuais (consistência eventual de R3/R4;
durabilidade da fila in-process de R4) foram avaliados como aceitáveis para
onboarding corporativo e documentados como DTA residual --- substituível por
BullMQ + Redis quando o SLA de durabilidade justificar o custo operacional.

\section{Resultados}
%% ─────────────────────────────────────────────────────────────────────────────

\subsection{Resultados dos testes de carga}

O Cenário~A apresentou ruptura operacional a 24\% da capacidade-alvo: o ponto de
degradação crítica ocorreu em 2.413 VUs simultâneos, com taxa de erro de 61,94\% e
latência média $>$\,5.000\,ms. A causa raiz foi o padrão N+1 na rota
\texttt{GET /app/team}, com até 109 queries síncronas por requisição e pool
limitado a 10 conexões, o banco esgotou conexões em dezenas de requisições
simultâneas. O Cenário~B, em contraste, completou todos os cinco estágios do
protocolo sem atingir ponto de ruptura\footnote{Momento em que o sistema deixa de
conseguir responder às requisições dentro dos critérios de SLA, caracterizado por
taxa de erro ou latência acima dos limiares definidos.}: a taxa de erro permaneceu em
0,00\% em todos os estágios, com latência média de 34\,ms e
throughput\footnote{Volume de requisições processadas com sucesso pelo servidor por
segundo; mede a capacidade de atendimento do sistema sob carga.} médio de 161\,req/s,
mesmo com o protocolo estendido até 10.000 VUs, o dobro do SLA definido.
A Tabela~\ref{tab:resultados} consolida as métricas dinâmicas e estáticas dos dois
cenários.

\begin{table}[h]
  \caption{Resultados comparativos: métricas de carga e estáticas, Cenário~A vs.\ Cenário~B.}
  \label{tab:resultados}
  \begin{tabular}{lrrl}
    \toprule
    \textbf{Métrica} & \textbf{Cen.~A} & \textbf{Cen.~B} & \textbf{SLA / ref.} \\
    \midrule
    \multicolumn{4}{l}{\textit{Dinâmicas (teste de carga)}} \\
    Ruptura (VUs)      & 2.413        & $>$\,10.000      & $\geq$\,5.000 \\
    Taxa de erro (\%)  & 61,94        & 0,00             & $<$\,1        \\
    Latência (ms)      & $>$\,5.000   & 34               & $<$\,500      \\
    Throughput (req/s) & ---          & 161              & ---           \\
    Queries/req        & 109 (N+1)    & 1 (JOIN)         & $O(1)$        \\
    \midrule
    \multicolumn{4}{l}{\textit{Estáticas (código-fonte)}} \\
    LOC\footnote{\textit{Lines of Code}: número de linhas de código-fonte.}
                       & $\approx$380 & 474              & $\Delta$\,+25\% \\
    CC média/rota      & 2,8          & 3,2              & $\Delta$\,+14\% \\
    Loops N+1          & 2            & 0                & $-$100\%      \\
    Pool máx.\ conexões & 10, s/\,timeout & 20, t/o\,2\,s & $\Delta$\,+100\% \\
    Cache              & Ausente      & TTL 60--300\,s   & ---           \\
    Escrita assíncrona & Síncrona     & Fila 500\,ms     & ---           \\
    \bottomrule
  \end{tabular}
\end{table}

A CC\footnote{Complexidade Ciclomática: número de caminhos independentes de execução
por módulo.} média do Cenário~B (3,2/rota) é ligeiramente superior à do Cenário~A
(2,8/rota). Esse aumento é esperado e desejável: resulta das verificações de cache,
do handler\footnote{Bloco de código responsável por tratar um evento ou erro
específico; no Cenário~B, há handlers para erros de pool esgotado e para timeouts de
requisição.} da fila e dos timeouts\footnote{Tempo máximo que o sistema aguarda por
uma resposta antes de encerrar a operação e retornar erro ao usuário.} explícitos.
O dado crítico é a eliminação dos dois loops \texttt{for/await}\footnote{Estrutura
de repetição em JavaScript que executa operações assíncronas sequencialmente; no
Cenário~A, era usada para consultar o banco repetidamente para cada colaborador,
gerando o padrão N+1.}: sua ausência transforma o comportamento de escala de
$O(N)$\footnote{Número de operações cresce proporcionalmente ao tamanho da entrada;
no Cenário~A, cada colaborador adicional gerava 2 consultas extras ao banco.} para
$O(1)$\footnote{Número de operações permanece constante independentemente da entrada;
no Cenário~B, uma única consulta com JOIN retorna todos os colaboradores de uma vez.}
queries por requisição, confirmando estaticamente a causa do colapso observado no
Cenário~A.

%% ─────────────────────────────────────────────────────────────────────────────
\section{Considerações finais}
%% ─────────────────────────────────────────────────────────────────────────────

\subsection{Contribuições e implicações práticas}

A principal contribuição empírica é a quantificação experimental da DTA gerada
por IAG em condições reais de startup. O ponto de ruptura em 2.413 VUs, com SLA
definido em 5.000, estabelece uma métrica concreta para a discussão da viabilidade
de MVPs gerados por IAG em contextos de escala. Em contraste, o Cenário~B não
apresentou ruptura até 10.000 VUs (0,00\% de erro, 34\,ms de latência), evidenciando
que as quatro refatorações arquiteturais eliminaram a causa raiz do colapso: o
padrão N+1 aliado ao pool subdimensionado.

A contribuição científica original é o \textbf{modelo de prompt arquitetural de
quatro componentes}. Enquanto a literatura de \textit{prompt engineering} propõe
padrões genéricos~\cite{White2023} e a literatura de RE documenta a lacuna de
requisitos não-funcionais em código gerado por IAG~\cite{Nascimento2023,Vogelsang2024},
este trabalho propõe e valida empiricamente um padrão específico para escalabilidade
arquitetural, operacionalizável por founders sem formação técnica. O modelo preenche
uma lacuna na interseção de três campos: (i)~\textit{prompt engineering} estruturado
para software; (ii)~elicitação de requisitos não-funcionais mediada por linguagem
natural; e (iii)~arquitetura de software assistida por IAG. Para founders, o estudo
demonstra que o caminho do Cenário~A ao~B é percorrível sem abandonar as ferramentas
de IAG: a lacuna identificada é de especificação de requisitos, não de capacidade
das ferramentas.

\subsection{Limitações}
\label{sec:limitacoes}

Quatro limitações: (1)~dados de negócio sintéticos; (2)~variabilidade do ambiente de
nuvem compartilhada (4 vCPUs, 16 GB RAM), que limita a reprodutibilidade exata;
(3)~fundamento em uma única startup, restringindo a generalização dos padrões P1--P5;
e (4)~avaliação qualitativa-comparativa realizada por um único autor, sem rubrica
formalizada.

\subsection{Trabalhos futuros}

Duas direções promissoras: (1)~extensão do protocolo de entrevistas para
identificar P6--P10 em startups de diferentes setores e tecnologias; e
(2)~desenvolvimento de um catálogo de prompts arquiteturais --- uma linguagem de
especificação de escalabilidade para uso com IAG --- como contribuição direta à
comunidade de Engenharia de Software.

%% ─────────────────────────────────────────────────────────────────────────────
\section*{Disponibilidade de artefatos}
%% ─────────────────────────────────────────────────────────────────────────────

Os artefatos deste estudo (scripts de teste de carga, Dockerfiles, prompts
utilizados nos Cenários~A e B e dados de métricas estáticas) estão disponíveis em
repositório público\footnote{\url{https://github.com/anneestherlf/welcome-navigator}}~\cite{WelcomeNavigator2025}. Os
dados qualitativos da entrevista não podem ser disponibilizados publicamente em virtude
do compromisso de anonimização firmado com os participantes. Os fundadores consentiram
com o uso dos dados para fins acadêmicos, com a condição de que nenhuma informação
identificadora fosse divulgada.

%% ─────────────────────────────────────────────────────────────────────────────
\section*{Agradecimentos}
%% ─────────────────────────────────────────────────────────────────────────────

As autoras agradecem aos fundadores da startup que gentilmente participaram da
entrevista e consentiram com o uso dos dados para fins acadêmicos.

Em conformidade com a política de uso de Inteligência Artificial Generativa do
VII~CTIC-ES, declara-se que o modelo Claude (Anthropic) foi utilizado como ferramenta
de auxílio na produção deste artigo, especificamente para: geração do código
\LaTeX{} das tabelas e da Figura~1; sugestão de reformulações textuais para adequação
ao registro científico; e geração da estrutura inicial de partes do texto, revisada e
aprovada pelas autoras. Todo o conteúdo científico (hipóteses, metodologia,
implementação dos cenários, execução dos testes e interpretação dos resultados)
é de autoria exclusiva das pesquisadoras.

%% ─────────────────────────────────────────────────────────────────────────────
\bibliographystyle{ACM-Reference-Format}
\bibliography{referencias}
%% ─────────────────────────────────────────────────────────────────────────────

\end{document}
